\documentclass[11pt,a4paper]{article}
\usepackage{auditstyle}
\lhead{Fresh Glimm--Jaffe proof audit}
\rhead{16 August 2026}
\usepackage{bm}
\newcommand{\eps}{\varepsilon}
\newcommand{\T}{\mathbb T}
\newcommand{\Z}{\mathbb Z}
\newcommand{\N}{\mathbb N}
\newcommand{\PP}{\mathbb P}
\newcommand{\U}{\mathcal U}
\newcommand{\D}{\mathcal D}
\newcommand{\HH}{\mathcal H}
\newcommand{\Tr}{\operatorname{Tr}}
\newcommand{\Jone}{\mathcal J_1}
\newcommand{\Jtwo}{\mathcal J_2}
\DeclareMathOperator{\cth}{coth}

\title{Part II V3 in Lamport notation\\[3pt]
\large Construction of the continuum torus Hamiltonian, trace class, periodic
Feynman--Kac formula, and the gap}
\author{Fresh reconstruction of \texttt{part\_ii/v3\_continuum\_package.tex}}
\date{16 August 2026}

\begin{document}
\maketitle

\begin{resultbox}
The continuum package is reconstructed below without importing a pre-existing
$P(\varphi)_2$ Hamiltonian.  The only printed input is Simon's four-result
second-quantization cluster, checked in the accompanying citation-evidence dossier.
The periodic cutoff estimate that V3 states only by analogy is supplied explicitly
in Steps $\langle3\rangle8$--$\langle3\rangle10$.
\end{resultbox}
\tableofcontents

\section{Data, exact target, and imports}

Fix $L,m>0$, $\lambda\ge0$, and $a,b\in\R$.  Put
\[
 \Gamma_\infty=\frac\pi L\Z,
 \qquad \gamma(k)=(m^2+k^2)^{1/2},
 \qquad \T_L=\R/(2L\Z).
\tag{0.1}
\]
Whenever a periodic time interval is used, fix $t\in[t_0,T]$ and impose
\[
 t_0L\ge1.\tag{0.2}
\]
The gate ultimately uses $(a,b)=(0,0)$.  Allowing arbitrary $(a,b)$ here makes
explicit that no continuum counterterm is being hidden.

Let $(e_j)_{j\ge1}$ be the real orthonormal Fourier basis
\[
 e_0=(2L)^{-1/2},\qquad e_k^c=L^{-1/2}\cos(kx),\qquad
 e_k^s=L^{-1/2}\sin(kx)\quad(k>0),
\tag{0.3}
\]
and let $\gamma_j$ be the corresponding value of $\gamma$.  On
\[
 (Q,\mu)=\left(\R^\N,\bigotimes_{j\ge1}\mathcal N(0,1)\right)
\tag{0.4}
\]
write the coordinate variables as $\xi_j$ and define
\[
 \varphi^{(K)}(x)=\sum_{|k_j|\le K}
 \frac{\xi_j}{\sqrt{2\gamma_j}}e_j(x),
 \qquad
 c_K=\E_\mu[\varphi^{(K)}(x)^2]
 =\frac1{2L}\sum_{|k|\le K}\frac1{2\gamma(k)}.
\tag{0.5}
\]
The free Hamiltonian on $L^2(Q,\mu)$ is
\[
 H_0=d\Gamma(\gamma),\qquad e^{-sH_0}=\Gamma(e^{-s\gamma}),
 \qquad
 Z_0(s)=\Tr e^{-sH_0}=\prod_{k\in\Gamma_\infty}(1-e^{-s\gamma(k)})^{-1}.
\tag{0.6}
\]
Its bounded smooth cylinder core is
\[
 \D=\{F(\xi_1,\ldots,\xi_n):n<\infty, F\in C_b^\infty(\R^n)\}.
\tag{0.7}
\]

The following are the only external statements used.

\begin{assumption}[Simon second-quantization cluster]
For a real-Hilbert-space contraction $A$, the Gaussian second quantization
$\Gamma(A)$ is positivity preserving; it is an $L^p$ contraction for every
$1\le p\le\infty$; it is positivity improving if $\|A\|<1$; and
\[
 \|\Gamma(A)f\|_q\le\|f\|_p
 \quad\hbox{if}\quad
 1<p\le q<\infty,\qquad
 \|A\|\le\sqrt{\frac{p-1}{q-1}}.
\tag{0.8}
\]
These are Simon, Theorem I.12, Corollary I.15, Theorems I.16 and I.17.
The exact pages are reproduced in the evidence dossier.
\end{assumption}

The elementary bounded-perturbation Trotter product formula and the spectral
theorem are used in the forms proved in the model-sheet reconstruction.  All
remaining inputs below are results already proved in the V1 and repaired V2
Lamport modules.

\section{Interaction and sharp-time convergence}

\LS{1}{1}{Bound the cutoff variance.}

\LS{2}{1}{For $K\ge m$,}
\[
 c_K\le C_c(m,L)\,[1+\log(1+K/m)].\tag{1.1}
\]

\Because Pair the sine and cosine modes and compare the decreasing summand with
the integral:
\[
 \sum_{0<k\le K}\frac1{\sqrt{m^2+k^2}}
 \le \frac L\pi\int_0^K\frac{du}{\sqrt{m^2+u^2}}+\frac1m
 =\frac L\pi\operatorname{arsinh}(K/m)+\frac1m.
\]
Since $\operatorname{arsinh}x\le1+\log(1+x)$, (1.1) follows.

\LS{2}{2}{This proves $\langle1\rangle1$.}\QedStep

\LS{1}{2}{Define and bound the mode-cutoff interaction.}

For $K\in\Gamma_\infty$, set
\[
 V^{(K)}=\lambda\int_{\T_L}
 \left[{:}\varphi^{(K)}(x)^4{:}_{c_K}
 +a{:}\varphi^{(K)}(x)^2{:}_{c_K}+b\right]dx.
\tag{1.2}
\]

\LS{2}{1}{For $u\in\R$ and $c\ge0$,}
\[
 {:}u^4{:}_c+a{:}u^2{:}_c+b
 =(u^2-3c)^2+a(u^2-3c)-6c^2+2ac+b.
\tag{1.3}
\]

\Because Expanding the right side gives
$u^4+(a-6c)u^2+3c^2-ac+b$, which is
$(u^4-6cu^2+3c^2)+a(u^2-c)+b$.

\LS{2}{2}{Consequently,}
\[
 {:}u^4{:}_c+a{:}u^2{:}_c+b
 \ge-\left(6c^2+2|a|c+\frac{a^2}{4}+|b|\right).
\tag{1.4}
\]

\Because For $v=u^2-3c$, $v^2+av=(v+a/2)^2-a^2/4\ge-a^2/4$;
$2ac+b\ge-2|a|c-|b|$.

\LS{2}{3}{Thus}
\[
 V^{(K)}\ge-b_K,\qquad
 b_K:=2L\lambda\left(6c_K^2+2|a|c_K+\frac{a^2}{4}+|b|\right)
 \le C[1+\log(1+K/m)]^2.
\tag{1.5}
\]

\LS{2}{4}{If $\lambda>0$, $V^{(K)}(\xi)\to+\infty$ as
$|\xi|\to\infty$ in the finitely many retained coordinates.}

\Because Cauchy--Schwarz on a space of volume $2L$ and Parseval give
\[
 \int_{\T_L}(\varphi^{(K)})^4dx
 \ge\frac1{2L}\left(\int_{\T_L}(\varphi^{(K)})^2dx\right)^2,
 \qquad
 \int_{\T_L}(\varphi^{(K)})^2dx
 =\sum_{|k_j|\le K}\frac{\xi_j^2}{2\gamma_j}.
\tag{1.6}
\]
All Wick-subtraction and $a$ terms are affine in the last sum.  A positive
quadratic polynomial in that sum therefore dominates them at infinity.  When
$\lambda=0$, the free harmonic operator supplies coercivity.

\LS{2}{5}{This proves $\langle1\rangle2$.}\QedStep

\LS{1}{3}{Construct the limiting sharp-time interaction in $L^2(\mu)$.}

\LS{2}{1}{Own-variance Wick ordering gives the exact chaos decomposition}
\[
 V^{(K)}=\lambda I_4(f_K)+\lambda a I_2(g_K)+2L\lambda b,
\tag{1.7}
\]
where $f_K,g_K$ are restrictions of fixed infinite-mode kernels.

\Because ${:}X^r{:}_{\E X^2}$ is the $r$th homogeneous Wiener chaos.
The integral imposes spatial momentum conservation and introduces no lower-chaos
terms because the Wick constant in (1.2) is exactly the variance in (0.5).

\LS{2}{2}{For $K'\ge K$, the quartic tail satisfies}
\[
 \left\|I_4(f_{K'})-I_4(f_K)\right\|_2^2
 =4!(2L)^{-2}
 \sum_{\substack{k_1+\cdots+k_4=0\\
       \max_i|k_i|\in(K,K']}}
 \prod_{i=1}^4\frac1{2\gamma(k_i)}.
\tag{1.8}
\]

\Because The Wick isometry contracts only equal chaos orders.  Expanding four
Fourier sums and integrating $x,y\in\T_L$ supplies two Kronecker constraints;
one is eliminated, leaving (1.8), including the factor $4!$.

\LS{2}{3}{Let $g(k)=(m+|k|)^{-1}$.  Then $(2\gamma(k))^{-1}\le g(k)$ and}
\[
 (g*g)(k)\le CL\frac{1+\log(2+|k|/m)}{m+|k|},\qquad
 (g*g*g)(k)\le CL^2\frac{[1+\log(2+|k|/m)]^2}{m+|k|}.
\tag{1.9}
\]

\Because $m+|k|\le2\sqrt{m^2+k^2}$.  For the first convolution split the
summation into $|u|\le|k|/2$, $|k-u|\le|k|/2$, and the complement.  In the
first region the second denominator is at least $(m+|k|)/2$ and
$\sum_{|u|\le R}(m+|u|)^{-1}\le CL[1+\log(1+R/m)]$; the second region is
symmetric.  In the complement both denominators are at least a fixed multiple of
$m+|k|$, and the remaining square tail is $O(L/(m+|k|))$.  Repeat the same
three-region split after convolving the first estimate with $g$; the partial sum of
$(1+\log)/(m+|u|)$ contributes the second logarithm.  These are finite sums on the
grid $(\pi/L)\Z$, so no continuum-limit interchange occurs.

\LS{2}{4}{The part of (1.8) with $|k_1|>K$ is at most}
\[
 CL^3\sum_{|k_1|>K}
 \frac{[1+\log(2+|k_1|/m)]^2}{(m+|k_1|)^2}
 \le CL^3\frac{[1+\log(1+K/m)]^2}{K/m}.
\tag{1.10}
\]

\Because Sum the remaining three momenta by (1.9).  Integral comparison of the
decreasing tail $\log^2 u/u^2$ gives the final bound.  The union over the four
choices of the large momentum changes only $C$.

\LS{2}{5}{The quadratic tail has variance}
\[
 \left\|I_2(g_{K'})-I_2(g_K)\right\|_2^2
 =\frac12\sum_{K<|k|\le K'}\gamma(k)^{-2}\le\frac{CL}{K}.
\tag{1.11}
\]

\Because In the real sine/cosine basis it is the sum of independent multiples of
$\xi_j^2-1$, whose variance is $2$.  The mode normalizations give the factor
$1/2$; compare the remaining sum with
$(2L/\pi)\int_K^\infty(m^2+u^2)^{-1}du$.

\LS{2}{6}{There is a chaos-$\le4$ random variable $V_L\in L^2(\mu)$ such that}
\[
 \|V^{(K)}-V^{(K')}\|_2+\|V^{(K)}-V_L\|_2
 \le C_V\frac{1+\log(1+K/m)}{(K/m)^{1/2}}.
\tag{1.12}
\]

\Because Combine (1.10), (1.11), orthogonality of distinct chaoses, and
completeness of $L^2(\mu)$.  The constant chaos cancels.  The limiting kernels
are precisely the truncation-convention objects denoted
$\lambda\int({:}\varphi^4{:}_{C_L}+a{:}\varphi^2{:}_{C_L}+b)dx$.

\LS{2}{7}{$V^{(K)}\to V_L$ in every finite $L^p(\mu)$.}

\Because If $X$ is a sum of chaoses of orders at most four, (0.8) applied with
$A=e^{-\tau}\mathbf1$, $q=1+e^{2\tau}$, and the orthogonal chaos decomposition
gives
\[
 \|X\|_q\le(q-1)^2\|X\|_2\qquad(q\ge2).
\tag{1.13}
\]
Use (1.13) for the difference and interpolate for $p<2$.

\LS{2}{8}{This proves $\langle1\rangle3$.}\QedStep

\section{The three Gaussian path ensembles and uniform exponential bounds}

\LS{1}{4}{Fix the three ensembles and their covariance identities.}

\LS{2}{1}{The vacuum ensemble is $(Q,\mu)$ of (0.4).}

\LS{2}{2}{The periodic ensemble $\mu^{\rm per}_{t,L}$ is the centered field on
$\T_t\times\T_L$ with per-mode covariance}
\[
 c_\omega(s-s')=
 \frac{\cosh(\omega(t/2-|s-s'|_{\T_t}))}{2\omega\sinh(\omega t/2)},
 \qquad \omega=\gamma(k).
\tag{2.1}
\]

\LS{2}{3}{The stationary Ornstein--Uhlenbeck ensemble is the Gaussian process
$X_s=(\xi_j(s))_j$ with}
\[
 \E[\xi_j(s)\xi_j(s')]=e^{-\gamma_j|s-s'|}.
\tag{2.2}
\]
For bounded cylinder $F,G$,
\[
 \E[\overline{F(X_0)}G(X_s)]
 =\ip{F}{\Gamma(e^{-s\gamma})G}_{L^2(\mu)}.
\tag{2.3}
\]

\Because Both sides of (2.3) are finite-dimensional Gaussian integrals with unit
marginals and cross covariance $e^{-s\gamma_j}$ in coordinate $j$.  Equality for
cylinders follows from their characteristic functions; $L^2$ density extends it.

\LS{2}{4}{This proves $\langle1\rangle4$.}\QedStep

\LS{1}{5}{Prove the abstract uniform Nelson lemma used at every cutoff.}

Let $A_K$ be chaos-$\le4$ random variables indexed by a full tail
$\mathcal K\subset\Gamma_\infty$.  Define
$\ell_K=1+\log(1+K/m)$ and assume
\begin{align}
 \ell_{K^+}-\ell_K&\le\log2,\tag{2.4}\\
 A_K&\ge-c_1\ell_K^2,\tag{2.5}\\
 \|A_K-A_{K'}\|_2&\le C_A\ell_K(K/m)^{-\kappa}\quad(K'\ge K),
 \qquad \kappa>0.\tag{2.6}
\end{align}

\LS{2}{1}{For $b$ with $L_b=\sqrt{b/(2c_1)}$ sufficiently large and
$b\le c_1\ell_K^2$, choose the largest $K_b\le K$ for which
$\ell_{K_b}\le L_b$.  Then}
\[
 \frac{L_b}{2}<\ell_{K_b}\le L_b.
\tag{2.7}
\]

\Because Maximality gives $\ell_{K_b^+}>L_b$ and (2.4) gives
$\ell_{K_b}>L_b-\log2\ge L_b/2$.  If $b>c_1\ell_K^2$, (2.5) makes
$\PP(A_K\le-b)=0$.

\LS{2}{2}{$A_{K_b}\ge-b/2$, hence}
\[
 \{A_K\le-b\}\subset
 \{A_K-A_{K_b}\le-b/2\}.
\tag{2.8}
\]

\LS{2}{3}{The $L^2$ scale of the difference obeys}
\[
 \rho_b:=\|A_K-A_{K_b}\|_2
 \le2C_AL_b e^{-\kappa(L_b/2-1)}.
\tag{2.9}
\]

\Because $K_b/m=e^{\ell_{K_b}-1}-1$.  For large $L_b$, this is at least
$\tfrac12e^{\ell_{K_b}-1}$; insert (2.7) into (2.6).

\LS{2}{4}{For every integer $q\ge1$,}
\[
 \PP(A_K\le-b)
 \le\left[\frac{2(2q-1)^2\rho_b}{b}\right]^{2q}.
\tag{2.10}
\]

\Because Apply Markov to $|A_K-A_{K_b}|^{2q}$ in (2.8) and apply
(1.13) with exponent $2q$ to the chaos-$\le4$ difference.

\LS{2}{5}{Choose $q$ so that $(2q)^2\le e^{\kappa L_b/4}<(2q+2)^2$.
After increasing the fixed threshold $b_1$,}
\[
 \PP(A_K\le-b)\le\exp[-e^{\kappa L_b/8}]
 \qquad(b\ge b_1),
\tag{2.11}
\]
uniformly in $K$.

\Because Substitute (2.9) into (2.10).  The logarithm of the bracket is
at most $-\kappa L_b/8$ after the polynomial factors in $L_b$ are absorbed;
$2q\ge\tfrac12e^{\kappa L_b/8}$ gives (2.11), after another harmless
increase of $b_1$.

\LS{2}{6}{For every $p<\infty$,}
\[
 \sup_K\E[e^{-pA_K}]<\infty.
\tag{2.12}
\]

\Because Layer cake gives
\[
 \E[e^{-pA_K}]
 \le e^{pb_1}+p\int_{b_1}^\infty e^{pb}\PP(A_K\le-b)\,db.
\]
Since $L_b$ is a positive constant times $\sqrt b$, the double-exponential
decay in (2.11) dominates $e^{pb}$ and makes the integral finite uniformly in
$K$.  If $A_K\to A$ in $L^2$, take an a.s. convergent subsequence and apply
Fatou to obtain the same bound for $e^{-A}$.

\LS{2}{7}{The bounded successor-gap hypothesis (2.4) holds on every full lattice
tail.}

\Because For $K^+=K+\pi/L$ and $K\ge\pi/L$,
\[
 \ell_{K^+}-\ell_K
 =\log\left(1+\frac{\pi/L}{m+K}\right)<\log2.
\]
This condition is essential to the scale selection (2.7); (2.5)--(2.6) alone do
not imply that selection for a tower-sparse index set.

\LS{2}{8}{This proves $\langle1\rangle5$.}\QedStep

\LS{1}{6}{Instantiate the Nelson lemma in all three ensembles.}

\LS{2}{1}{For $0<\sigma\le\sigma_0$,}
\[
 \sup_K\|e^{-\sigma V^{(K)}}\|_{L^p(\mu)}<\infty,\qquad
 e^{-\sigma V_L}\in L^p(\mu).
\tag{2.13}
\]

\Because Conditions (2.5) and (2.6) are respectively (1.5) and (1.12),
with $\kappa=1/2$, after multiplying constants by $\sigma_0$.

\LS{2}{2}{For}
\[
 \mathcal W_K(t)=\int_0^tV^{(K)}(X_s)ds,\qquad
 \mathcal W(t)=\int_0^tV_L(X_s)ds,
\tag{2.14}
\]
one has, uniformly for $0<t\le T$,
\[
 \|\mathcal W_K-\mathcal W_{K'}\|_2
 \le t\|V^{(K)}-V^{(K')}\|_2,\quad
 \mathcal W_K\ge-TC\ell_K^2,\quad
 \sup_K\|e^{-\mathcal W_K}\|_p<\infty.
\tag{2.15}
\]

\Because Stationarity and Cauchy--Schwarz give
\[
 \|\mathcal W_K-\mathcal W_{K'}\|_2^2
 =\int_0^t\!\int_0^t\E[\Delta V(X_s)\Delta V(X_u)]dsdu
 \le t^2\|\Delta V\|_2^2.
\]
The pathwise lower bound is the integral of (1.5).  Time integration of a
global Gaussian chaos remains in chaos orders at most four, so Step
$\langle1\rangle5$ applies.

\LS{2}{3}{Let $\U_t^{(K)}$ be the spatial-mode truncation of the periodic
interaction.  Then}
\[
 \U_t^{(K)}\to\U_t\quad\hbox{in }L^2(\mu^{\rm per}_{t,L}),\qquad
 \sup_K\|e^{-\U_t^{(K)}}\|_p<\infty.
\tag{2.16}
\]

\LS{3}{1}{The pathwise lower bound is $\U_t^{(K)}\ge-Ct\ell_K^2$.}

\Because Apply (1.4) at every spacetime point using the vacuum cutoff constant
$c_K$, which obeys (1.1), and integrate over $\T_t\times\T_L$.

\LS{3}{2}{The quadratic chaos tail satisfies}
\[
 \|Y_2^{(K')}-Y_2^{(K)}\|_2^2
 \le C\sum_{|k|>K}\sum_{k_0\in(2\pi/t)\Z}
 (m^2+k_0^2+k^2)^{-2}\le C(t,L,m)K^{-2}.
\tag{2.17}
\]

\Because The first inequality is the Wick isometry and the continuum covariance
envelope.  Summing $k_0$ by comparison with its integral gives
$\sum_{k_0}(A^2+k_0^2)^{-2}\le C(t_0)(A^{-4}+tA^{-3})$; the remaining
one-dimensional spatial tail is $O(tLK^{-2})$ (and the $A^{-4}$ term is
smaller under $tL\ge1$).

\LS{3}{3}{The quartic chaos tail satisfies the explicit bound}
\[
 \|Y_4^{(K')}-Y_4^{(K)}\|_2^2
 \le C(t,L,m)
 \sum_{\substack{|k|>K\\ k_0\in(2\pi/t)\Z}}
 \frac{[1+\log(2+|(k_0,k)|)]^2}
 {(m^2+k_0^2+k^2)^2}
 \le C(t,L,m)K^{-2}[1+\log K]^2.
\tag{2.18}
\]

\Because Select one of the four momenta whose spatial component exceeds $K$;
there are four choices.  Sum the other three constrained covariance factors using
the two-dimensional convolution estimate already proved in the repaired V2 module:
$c_{\rm env}^{*3}(-\kappa)\le C(tL)^2[1+\log(2+|\kappa|)]^2/
(m^2+|\kappa|^2)$.  Multiplication by the selected factor gives the first sum.
Cover the anisotropic grid by its fundamental rectangles.  Its radial majorant is
decreasing outside a fixed disk, so the sum is bounded by its density times
\[
 \int_{r>K}\frac{r[1+\log(2+r)]^2}{(m^2+r^2)^2}dr
 \le C\frac{[1+\log K]^2}{m^2+K^2}.
\]
This is the missing ``same tail counts'' calculation in the source V3 document.

\LS{3}{4}{The thermal recombination coefficient tail is exponentially small.}

\Because
\[
 d^{(K)}(t)=\frac1{2L}\sum_{|k|\le K}
 \frac1{\gamma(k)(e^{t\gamma(k)}-1)},
\]
and for $t\ge t_0$ the omitted summand is bounded by
$C e^{-t_0|k|}/(m+|k|)$.  Hence
$|d^{(K)}-d^{(K')}|\le Ce^{-t_0K/2}$ after enlarging $C$.

\LS{3}{5}{Equations (2.17)--(2.18), the finite Wick recombination identity, and
boundedness of $d^{(K)}$ imply}
\[
 \|\U_t^{(K)}-\U_t^{(K')}\|_2
 \le C[1+\log(1+K/m)]^2(K/m)^{-1/2}.
\tag{2.19}
\]

\Because The actual chaos tails from (2.17)--(2.18) are at least as small as
$K^{-1}\log K$ after taking square roots.  The displayed $K^{-1/2}\log^2K$
rate is a weaker common bound that also covers the coefficient cross terms.  Each
cross term is a bounded coefficient difference times a bounded lower-chaos norm.

\LS{3}{6}{Absorb one logarithm by
$\ell_K\le C_\varsigma(K/m)^\varsigma$, choose
$0<\varsigma<1/2$, and apply Step $\langle1\rangle5$ with
$\kappa=1/2-\varsigma$.  This proves (2.16).}

\LS{2}{4}{The deterministically shifted periodic interactions obey the same
Cauchy and exponential-integrability conclusions, uniformly for the bounded
coarse-data class used in V4.}

\Because Substitute $u=\eta^{(K)}(s,x)+H^{(K)}(s,x)$ in (1.4).  Since (1.4)
holds for every real $u$, integration gives the same lower bound
$\U_t^{(K)}(\eta+H^{(K)})\ge-Ct\ell_K^2$.  Expand each shifted Wick monomial by
the Appell identity (3.12) and its degree-two analogue.  The repaired V2 theorem,
whose fundamental-band and alias sums are displayed in its Steps
$\langle1\rangle7$--$\langle1\rangle8$, bounds the resulting random-chaos tails
and deterministic coefficient differences by the right side of (2.19), uniformly
over the stated coarse-data bounds.  These are precisely the lower-tail and
$L^2$-increment hypotheses of Step $\langle1\rangle5$; applying that step proves
the asserted uniform exponential integrability and $L^2$ convergence.

\LS{2}{5}{This proves $\langle1\rangle6$.}\QedStep

\LS{1}{7}{Derive uniform spectral lower bounds for the finite-cutoff operators.}

Let $K^{(K)}$ be defined in the next section.  Its finite-mode periodic trace
identity will be proved independently of this step.

\LS{2}{1}{For a fixed admissible $t_0$,}
\[
 e^{-t_0\inf\spec K^{(K)}}\le\Tr e^{-t_0K^{(K)}}
 =Z_0(t_0)\E[e^{-\U_{t_0}^{(K)}}]\le Z_0(t_0)B.
\tag{2.20}
\]

\Because The first inequality selects the lowest heat eigenvalue; the identity is
proved in Step $\langle1\rangle8$ below; the last inequality is (2.16) at $p=1$.
No uniform lower bound is used in proving that finite-mode identity, so this is not
circular.

\LS{2}{2}{Taking logarithms gives a $K$-independent constant $c_\lambda$ with}
\[
 K^{(K)}\ge-c_\lambda.
\tag{2.21}
\]

\LS{2}{3}{For $0<\delta<1$, apply (2.21) to the amplified coupling
$\lambda/(1-\delta)$ to obtain}
\[
 K^{(K)}=\delta H_0+(1-\delta)
 \left[H_0+(1-\delta)^{-1}V^{(K)}\right]
 \ge\delta H_0-c_\delta
\tag{2.22}
\]
as quadratic forms.

\LS{2}{4}{This proves $\langle1\rangle7$.}\QedStep

\section{Finite-mode systems and their exact Feynman--Kac formulas}

For fixed $K$, decompose
$L^2(Q,\mu)=L^2(\mu_{\le K})\otimes L^2(\mu_{>K})$ and set
\[
 K^{(K)}=(H_0^{\le K}+V^{(K)})\otimes\mathbf1
 +\mathbf1\otimes H_0^{>K},\qquad T_K(s)=e^{-sK^{(K)}}.
\tag{3.1}
\]

\LS{1}{8}{Construct each finite-mode system and prove every representation used
later.}

\LS{2}{1}{The form of $H_0^{\le K}+V^{(K)}$ is densely defined,
lower bounded, and closed.}

\Because Choose $C>b_K$ and write
\[
 q_K(f)=\sum_{|k_j|\le K}\gamma_j\|\partial_jf\|_2^2
 +\int(V^{(K)}+C)|f|^2d\mu_{\le K}.
\tag{3.2}
\]
Both summands are nonnegative and their domains contain
$C_c^\infty(\R^{n_K})$.  If $f_n$ is Cauchy in the norm
$\|f\|_2^2+q_K(f)$, it converges in $L^2$ and in each weighted weak derivative.
A subsequence converges pointwise a.e.; Fatou for the multiplication term and
closedness of weak derivatives identify the limits.  Thus the form norm is
complete.

\LS{2}{2}{The finite-mode operator has compact resolvent.}

\Because On every Euclidean ball the Gaussian density is bounded above and below
by positive constants, so the weighted form norm controls the ordinary $H^1$ norm
there; Rellich makes the restricted embedding compact.  Outside a sufficiently
large ball, coercivity from (1.6) gives
$(V^{(K)}+C)\ge R$, hence
$\int_{|\xi|>R_0}|f|^2d\mu\le R^{-1}q_K(f)$.  First take $R$ large and then
use Rellich on the ball.  The unit ball of the form domain is therefore relatively
compact in $L^2$.

\LS{2}{3}{For one retained coordinate, the free Mehler kernel with respect to
$\mathcal N(0,1)$ is}
\[
 M_s^{(j)}(x,y)=(1-r_j^2)^{-1/2}
 \exp\left[-\frac{r_j^2(x^2+y^2)-2r_jxy}{2(1-r_j^2)}\right],
 \qquad r_j=e^{-s\gamma_j}.
\tag{3.3}
\]

\Because Completing the square shows
$\int M_s^{(j)}(x,y)d\mu(y)=1$; symmetry is evident.  Applying the Gaussian
generating function to Hermite polynomials gives eigenvalue $r_j^n$ in degree
$n$.  Hence the product over retained coordinates is exactly
$\Gamma(e^{-s\gamma})$ on that finite Gaussian space.

\LS{2}{4}{For bounded $f,g$ on the retained coordinates,}
\[
 \ip f{e^{-s(H_0^{\le K}+V^{(K)})}g}
 =\E\left[\overline{f(X_0^{\le K})}g(X_s^{\le K})
 e^{-\int_0^sV^{(K)}(X_u)du}\right].
\tag{3.4}
\]

\LS{3}{1}{First replace $V^{(K)}$ by $V_R=V^{(K)}\wedge R$, with $R>0$.}

\Because By (1.5), $-b_K\le V_R\le R$, so multiplication by $V_R$ is
bounded and self-adjoint.

\LS{3}{2}{The bounded-perturbation Trotter formula gives}
\[
 e^{-s(H_0+V_R)}=s\!\!\!\lim_{n\to\infty}
 (e^{-(s/n)H_0}e^{-(s/n)V_R})^n.
\tag{3.5}
\]

\LS{3}{3}{Insertion of (3.3) between the multiplication factors turns the matrix
element of the product into the OU chain integral with weight}
\[
 \exp\left[-\frac{s}{n}\sum_{j=1}^nV_R(X_{js/n})\right].
\]

\LS{3}{4}{As $n\to\infty$ this weight converges a.s. to
$e^{-\int_0^sV_R(X_u)du}$ and is bounded by $e^{sb_K}$.}

\Because The finite-dimensional OU process has continuous paths and $V_R$ is
bounded continuous.  The displayed sum is therefore an ordinary Riemann sum;
dominated convergence applies.

\LS{3}{5}{As $R\to\infty$, the operators and path integrals converge to (3.4).}

\Because $V_R\uparrow V^{(K)}$, so the associated closed forms increase and
monotone form convergence gives strong semigroup convergence.  On each continuous
path, $e^{-\int V_R}\downarrow e^{-\int V^{(K)}}$; it is bounded by $e^{sb_K}$,
so dominated convergence applies.  This proves (3.4).

\LS{2}{5}{The kernel in (3.4) is jointly continuous and strictly positive.}

\Because Condition the OU law on its two endpoints.  The product Mehler density is
continuous and strictly positive.  The bridge expectation of the weight lies in
$(0,e^{sb_K}]$.  Continuity follows by dominated convergence after realizing the
finite-dimensional Gaussian bridge as an affine continuous function of its endpoints
plus an endpoint-independent bridge.

\LS{2}{6}{Tensoring with the independent free tail gives, for bounded cylinders
$f,g$,}
\[
 \ip f{T_K(s)g}
 =\E\left[\overline{f(X_0)}g(X_s)e^{-\mathcal W_K(s)}\right].
\tag{3.6}
\]

\Because The retained and omitted OU coordinates are independent; apply (3.4) to
the first factor and (2.3) to the second.

\LS{2}{7}{$T_K(s)$ is trace class and}
\[
 \Tr T_K(s)\le e^{sb_K}Z_0(s).
\tag{3.7}
\]

\Because The min--max principle and $V^{(K)}\ge-b_K$ imply that every ordered
eigenvalue of the retained interacting operator is at least the corresponding free
eigenvalue minus $b_K$.  Tensoring with the free tail and summing all occupation
numbers gives (3.7).  In particular, the full tensor operator has compact resolvent.

\LS{2}{8}{The periodic trace formula is}
\[
 \Tr T_K(s)=Z_0(s)\E_{\mu^{\rm per}_{s,L}}[e^{-\U_s^{(K)}}].
\tag{3.8}
\]

\Because Apply the diagonal trace to the strictly positive kernel from (3.4),
normalize the product of cyclic Mehler kernels by its free trace, and integrate the
intermediate points.  The normalized cyclic Gaussian chain has covariance (2.1):
for one mode, direct inversion of the cyclic tridiagonal quadratic form gives
$\cosh(\omega(s/2-|u-v|))/(2\omega\sinh(\omega s/2))$.  The equal-time
variance differs from $c_K$ by $d^{(K)}(s)$; the exact Wick-recombination identity
\[
 {:}X^4{:}_{c_K}={:}X^4{:}_{v_K}
 +6(v_K-c_K){:}X^2{:}_{v_K}+3(v_K-c_K)^2
\tag{3.9}
\]
identifies the path action with $\U_s^{(K)}$.  Passing from equal cyclic
partitions to the continuous bridge uses the bounded-truncation and dominated
limits already justified in Steps $\langle3\rangle1$--$\langle3\rangle5$.

\LS{2}{9}{For a Weyl symbol whose position and momentum components are supported
in $|k|\le K$,}
\[
 \frac{\Tr(T_K(s)W_{\rm ct}(\xi))}{\Tr T_K(s)}
 =e^{-S^{(K)}(H_p)}
 \frac{\E[e^{i\Phi(J_q)}e^{-\U_s^{(K)}(\eta+H_p^{(K)})}]}
 {\E[e^{-\U_s^{(K)}]}}.
\tag{3.10}
\]
Here
\[
 S^{(K)}(H_p)=\frac1{8L}\sum_{|k|\le K}
 \gamma(k)\cth(s\gamma(k)/2)|b_\infty(k)|^2.
\tag{3.11}
\]

\Because Apply the V1 proof mode by mode: the Weyl operator translates the endpoint
by its momentum component and multiplies by its position phase; the unique solution
of $(-\partial_u^2+\gamma(k)^2)H=0$ off the marked slice with the prescribed jump
removes the endpoint translation; completing the Gaussian square produces (3.11)
and the sharp-time phase.  The Wick Appell identity
\[
 {:}(X+H)^4{:}_c
 ={:}X^4{:}_c+4H{:}X^3{:}_c+6H^2{:}X^2{:}_c+4H^3X+H^4
\tag{3.12}
\]
gives the shifted action without an extra $-6c(2HX+H^2)$ term.  To remove the
interaction cutoff, replace $V^{(K)}$ by $V_R=V^{(K)}\wedge R$ exactly as in
Steps $\langle3\rangle1$--$\langle3\rangle5$: the Weyl insertion is unitary, the
truncated path weight is bounded by $e^{sb_K}$, and dominated convergence first
passes the Trotter mesh and then $R\to\infty$.  All retained spaces are
finite-dimensional.  Tensoring with the omitted free modes produces the same free
factor in numerator and denominator, so that factor cancels.

\LS{2}{10}{The kernel is positivity improving, so the finite-mode lowest
eigenvalue is simple.}

\Because If $f,g\ge0$ are nonzero, the endpoint Mehler density is strictly
positive and the Feynman--Kac weight is strictly positive on every finite path;
therefore $\ip f{e^{-s(H_0^{\le K}+V^{(K)})}g}>0$.  Let $T$ denote this compact
self-adjoint integral operator and let $Th=\|T\|h$.  Its strictly positive kernel
gives $|Th(x)|\le T|h|(x)$; the variational bound
$\langle |h|,T|h|\rangle\le\|T\|\|h\|^2$ forces equality.  Equality in the
integral triangle inequality for a strictly positive kernel forces $h(y)$ to have
one common phase for almost every $y$.  Thus every top eigenvector is a scalar
multiple of one strictly positive vector, so the top eigenvalue is simple.
Tensoring with the free tail preserves a unique total ground vector because each
strictly positive tail frequency has the unique oscillator vacuum.

\LS{2}{11}{This proves $\langle1\rangle8$ and supplies the identity used in
(2.20).}\QedStep

\section{Strong semigroup limit and the continuum generator}

\LS{1}{9}{Prove that $T_K(s)$ converges strongly to one self-adjoint semigroup.}

\LS{2}{1}{For a bounded cylinder $\psi$ supported in modes at most $K_0$ and
$K,K'\ge K_0$,}
\[
 \|(T_K(s)-T_{K'}(s))\psi\|_2
 \le\|\psi\|_\infty
 \|e^{-\mathcal W_K(s)}-e^{-\mathcal W_{K'}(s)}\|_2.
\tag{4.1}
\]

\Because Subtract the two conditional-expectation versions of (3.6), condition on
$X_0$, square, integrate, and apply conditional Jensen.

\LS{2}{2}{For real $x,y$,}
\[
 |e^{-x}-e^{-y}|\le|x-y|(e^{-x}+e^{-y}).\tag{4.2}
\]

\Because Integrate the derivative of $e^{-[(1-u)x+uy]}$ over $u\in[0,1]$ and
bound every intermediate exponential by $e^{-x}+e^{-y}$.

\LS{2}{3}{Equations (1.13), (2.13)--(2.15), and H\"older give}
\[
 \|e^{-\mathcal W_K}-e^{-\mathcal W_{K'}}\|_2
 \le C_Ts\|V^{(K)}-V^{(K')}\|_2.
\tag{4.3}
\]

\Because Apply (4.2), then
$\|XY\|_2\le\|X\|_4\|Y\|_4$ with
$X=\mathcal W_K-\mathcal W_{K'}$ and
$Y=e^{-\mathcal W_K}+e^{-\mathcal W_{K'}}$.  The first norm is at most
$9s\|V^{(K)}-V^{(K')}\|_2$ by (1.13), (2.15); the second is uniformly
bounded by (2.13)--(2.15).

\LS{2}{4}{$T_K(s)$ is strongly Cauchy on $L^2(\mu)$, uniformly for $s$ in a
compact interval.}

\Because Equations (4.1), (4.3), and (1.12) prove this on the dense bounded
cylinder core.  The uniform operator bound
$\|T_K(s)\|\le e^{sc_\lambda}$ from (2.21) extends it to all $L^2$.
Call the strong limit $T(s)$.

\LS{2}{5}{$T(s)$ is a self-adjoint semigroup with
$\|T(s)\|\le e^{sc_\lambda}$.}

\Because Strong limits preserve the adjoint identity on matrix elements.
For the product, insert and subtract
$T_K(s)T_K(u)$ and use uniform boundedness to pass both factors strongly:
$T(s)T(u)=T(s+u)$.  The norm bound passes by lower semicontinuity.

\LS{2}{6}{$T(s)$ is strongly continuous at $s=0$.}

\Because For a bounded cylinder $\psi$ and fixed $K$ containing its modes,
\[
 \|T(s)\psi-\psi\|
 \le\|T(s)\psi-T_K(s)\psi\|+\|T_K(s)\psi-\psi\|.
\]
By (4.1)--(4.3), the first term is at most
$Cs\|\psi\|_\infty\sup_{K'\ge K}\|V^{(K')}-V^{(K)}\|_2$; the second
tends to zero for fixed $K$.  First let $s\downarrow0$, then $K\to\infty$.
Density and the uniform norm bound finish.

\LS{2}{7}{There is a unique self-adjoint $K_L\ge-c_\lambda$ such that}
\[
 T(s)=e^{-sK_L}.
\tag{4.4}
\]

\Because Apply the spectral-generator theorem for a strongly continuous
self-adjoint semigroup bounded by $e^{sc_\lambda}$.  The semigroup uniquely
determines its generator.

\LS{2}{8}{For bounded measurable $f,g$,}
\[
 \ip f{T(s)g}=\E[\overline{f(X_0)}g(X_s)e^{-\mathcal W(s)}].
\tag{4.5}
\]

\Because For bounded cylinders, pass $K\to\infty$ in (3.6) by (4.3) and
strong convergence.  Approximate bounded measurable functions by their conditional
expectations on the increasing mode sigma-fields.  Martingale convergence gives
a.e. and $L^2$ convergence; dominated convergence on the right uses
$e^{-\mathcal W}\in L^1$, while strong continuity of $T(s)$ handles the left.

\LS{2}{9}{The uniform form bound passes to the limit:}
\[
 K_L\ge\tfrac12H_0-c_{1/2}.
\tag{4.6}
\]

\LS{3}{1}{Choose $M>c_{1/2}\vee c_\lambda$.  From (2.22),}
\[
 (K^{(K)}+M)^{-1}\le
 (\tfrac12H_0+M-c_{1/2})^{-1}.
\tag{4.7}
\]

\Because For strictly positive self-adjoint $S$, spectral calculus gives
\[
 q_S(\psi)=\sup_\phi[2\operatorname{Re}\ip\psi\phi-ip\phi{S^{-1}\phi}],
\quad
 \ip\phi{S^{-1}\phi}=\sup_\psi[2\operatorname{Re}\ip\psi\phi-q_S(\psi)].
\tag{4.8}
\]
Thus form order and the reverse resolvent order are equivalent.

\LS{3}{2}{The resolvents converge strongly:}
\[
 (K^{(K)}+M)^{-1}	o(K_L+M)^{-1}.
\tag{4.9}
\]

\Because Use the Laplace representation
$S^{-1}=\int_0^\infty e^{-Ms}e^{-s(S-M)}ds$ with the semigroups
$T_K,T$; $M>c_\lambda$ supplies an integrable uniform majorant.

\LS{3}{3}{Pass to the limit in (4.7) on every vector and use (4.8) in the
reverse direction.  This gives (4.6) and
$Q(K_L)\subset Q(H_0)$.}

\LS{2}{10}{This proves $\langle1\rangle9$.}\QedStep

\LS{1}{10}{Identify the generator and its form on the cylinder core.}

\LS{2}{1}{For $\psi\in\D$,}
\[
 \frac{(1-T(s))\psi}{s}
 =\frac{(1-e^{-sH_0})\psi}{s}
 -\frac1s\E[\psi(X_s)(e^{-\mathcal W(s)}-1)\mid X_0].
\tag{4.10}
\]

\Because This is (4.5) written as a conditional expectation, with the free
conditional expectation subtracted and added.

\LS{2}{2}{The first term in (4.10) converges in $L^2$ to $H_0\psi$.}

\Because Let $d\mu_\psi(\lambda)=d\langle\psi,E_{H_0}(\lambda)\psi\rangle$.
The spectral theorem gives
\[
 \left\|\frac{1-e^{-sH_0}}s\psi-H_0\psi\right\|_2^2
 =\int_0^\infty
 \left|\frac{1-e^{-s\lambda}}s-\lambda\right|^2d\mu_\psi(\lambda).
\]
The integrand converges pointwise to zero and is bounded by $\lambda^2$, because
$0\le(1-e^{-s\lambda})/s\le\lambda$.  Since $\psi\in D(H_0)$,
$\int\lambda^2d\mu_\psi<\infty$; dominated convergence proves the claim.

\LS{2}{3}{Write $e^{-x}-1=-x+\rho(x)$, where}
\[
 |\rho(x)|\le\tfrac12x^2e^{x_-},\qquad x_-=\max(-x,0).
\tag{4.11}
\]

\Because Taylor's formula with integral remainder gives
$\rho(x)=x^2\int_0^1(1-u)e^{-ux}du$; bound $e^{-ux}\le e^{x_-}$.

\LS{2}{4}{The linear part satisfies}
\[
 \frac1s\E[\psi(X_s)\mathcal W(s)\mid X_0]
 =\frac1s\int_0^s e^{-uH_0}
 \bigl(V_L e^{-(s-u)H_0}\psi\bigr)du
 \longrightarrow V_L\psi
\tag{4.12}
\]
in $L^2$.

\Because The Markov property gives the identity.  Since $V_L\in L^4$ and
$\psi$ is bounded,
\[
 \|V_L(e^{-vH_0}\psi-\psi)\|_2
 \le\|V_L\|_4\|e^{-vH_0}\psi-\psi\|_4.
\]
For $\psi=F(\xi_1,\ldots,\xi_n)$, the Mehler formula and bounded derivatives of
$F$ give
\[
 \sup_{0\le v\le s}\|e^{-vH_0}\psi-\psi\|_4
 \le C_F\sum_{j\le n}(1-e^{-s\gamma_j})^{1/2}\to0.
\]
Also $e^{-uH_0}$ is strongly continuous on $V_L\psi\in L^2$.  Both errors are
uniform for $0\le u\le s$, proving (4.12).

\LS{2}{5}{The remainder in (4.11), divided by $s$, tends to zero in $L^2$.}

\Because Conditional Jensen, H\"older, (1.13), and (2.15) give
\[
 \frac1s\|\E[\psi(X_s)\rho(\mathcal W)\mid X_0]\|_2
 \le\frac{\|\psi\|_\infty}{2s}
 \|\mathcal W\|_8^2\|e^{\mathcal W_-}\|_4
 \le Cs.
\tag{4.13}
\]
Indeed $\|\mathcal W\|_8\le49s\|V_L\|_2$ and
$e^{\mathcal W_-}\le1+e^{-\mathcal W}$, whose $L^4$ norm is uniformly
bounded by Step $\langle1\rangle6$.

\LS{2}{6}{Therefore}
\[
 \D\subset D(K_L),\qquad K_L\psi=H_0\psi+V_L\psi
 \quad(\psi\in\D).
\tag{4.14}
\]

\Because Combine (4.10)--(4.13) with the generator characterization
$K_L\psi=\lim_{s\downarrow0}s^{-1}(1-e^{-sK_L})\psi$.

\LS{2}{7}{For $\psi\in\D$,}
\[
 q_{K_L}(\psi)=\ip\psi{H_0\psi}+\int V_L|\psi|^2d\mu.
\tag{4.15}
\]

\Because Spectral calculus gives
$s^{-1}\ip\psi{(1-e^{-sK_L})\psi}\to q_{K_L}(\psi)$ after shifting by the
lower spectral bound.  Equation (4.14) identifies the same limit with the
right side.  This proves equality on $\D$ only; no unproved claim that $\D$ is a
form core is made.

\LS{2}{8}{This proves $\langle1\rangle10$.}\QedStep

\section{Uniform smoothing, trace class, and trace-norm approximation}

\LS{1}{11}{Prove the cutoff-uniform $L^2\to L^4$ smoothing estimate.}

Put $s_4=(\log3)/m$.

\LS{2}{1}{By self-adjoint duality it suffices to prove}
\[
 \sup_K\|T_K(s_4)\|_{4/3\to2}<\infty.
\tag{5.1}
\]

\Because $\|T_K(s_4)\|_{2\to4}=\|T_K(s_4)\|_{4/3\to2}$.  For
$s>s_4$, compose with
$\|T_K(s-s_4)\|_{2\to2}\le e^{c_\lambda(s-s_4)}$.

\LS{2}{2}{Replace $W=V^{(K)}$ by $W_R=W\wedge R$ and use Trotter with
$\tau=s_4/n$.  It is enough to bound}
\[
 (e^{-\tau W_R}e^{-\tau H_0})^n:L^{4/3}\to L^2
\tag{5.2}
\]
uniformly in $n,K,R$.

\Because The bounded Trotter products converge strongly.  The form semigroups for
$W_R\uparrow W$ converge strongly.  An a.e. convergent subsequence and Fatou pass
any common $L^2$ bound through both limits.

\LS{2}{3}{For $j=0,\ldots,n$, define}
\[
 p_j=1+\frac13e^{j\tau m},\qquad
 P_j=1+(p_j-1)e^{2\tau m}.
\tag{5.3}
\]
Then $p_0=4/3$, $p_n=2$, and
\[
 \|e^{-\tau H_0}\|_{p_j\to P_j}\le1.
\tag{5.4}
\]

\Because $\|e^{-\tau\gamma}\|=e^{-\tau m}$ and
$e^{-\tau m}=\sqrt{(p_j-1)/(P_j-1)}$; apply (0.8).

\LS{2}{4}{Define $r_j$ by}
\[
 \frac1{p_{j+1}}=\frac1{r_j}+\frac1{P_j}.
\tag{5.5}
\]
Then multiplication gives
\[
 \|e^{-\tau W_R}u\|_{p_{j+1}}
 \le\|e^{-\tau W_R}\|_{r_j}\|u\|_{P_j}.
\tag{5.6}
\]

\LS{2}{5}{For $n\ge2$,}
\[
 \sum_{j=0}^{n-1}\frac1{r_j}\le1,\qquad
 r_j\tau\le\frac{12}{m}.
\tag{5.7}
\]

\Because From (5.3),
$P_j-p_{j+1}=\tfrac13e^{(j+1)\tau m}(e^{\tau m}-1)$.  Since
$p_{j+1},P_j\ge4/3$,
\[
 \sum_jr_j^{-1}
 \le\frac3{16}e^{\tau m}(e^{n\tau m}-1)
 =\frac38e^{\tau m}\le\frac{3\sqrt3}{8}<1.
\]
Also $p_{j+1}\le2$, $P_j\le1+\tfrac13e^{(j+2)\tau m}$, and
$\tau m\le e^{\tau m}-1$, whence
\[
 r_j\tau m
 =\frac{3\tau m,p_{j+1}P_j}
 {e^{(j+1)\tau m}(e^{\tau m}-1)}
 \le6+2e^{\tau m}<12.
\]

\LS{2}{6}{There is a constant $D$ independent of $j,n,K,R$ such that}
\[
 \|e^{-\tau W_R}\|_{r_j}\le D^{1/r_j}.
\tag{5.8}
\]

\Because Since $R>0$, $e^{-r_j\tau W_R}\le1+e^{-r_j\tau W}$.
Equation (5.7) restricts $r_j\tau$ to $(0,12/m]$, so the vacuum Nelson
bound (2.13), with $p=1$ and all $\sigma$ in that compact interval, makes
$\E e^{-r_j\tau W}$ uniformly bounded.

\LS{2}{7}{Multiplication of (5.4), (5.6), and (5.8) yields}
\[
 \|(e^{-\tau W_R}e^{-\tau H_0})^n\|_{4/3\to2}
 \le D^{\sum_j1/r_j}\le\max(1,D)=:M_0.
\tag{5.9}
\]

\LS{2}{8}{Consequently, uniformly in $K$,}
\[
 \|T_K(s)\|_{2\to4}\le M_0e^{c_\lambda(s-s_4)}\quad(s\ge s_4),
\tag{5.10}
\]
and the same holds for $T(s)$ by strong convergence and Fatou.

\LS{2}{9}{This proves $\langle1\rangle11$.}\QedStep

\LS{1}{12}{Upgrade strong convergence to operator-norm convergence at one time.}

Set $s_2=2s_4$ and $\Delta V=V^{(K')}-V^{(K)}$.

\LS{2}{1}{For bounded cylinder $\phi,\psi$ and every $0<s<s_2$, both
$T_K(s)\phi$ and $T_{K'}(s_2-s)\psi$ lie in the form domains of both
$K^{(K)}$ and $K^{(K')}$.}

\Because The uniform form lower bound (2.22) puts both vectors in $Q(H_0)$.
From (3.6), conditional Jensen, and the OU Nelson bound,
\[
 \|T_K(u)\phi\|_4\le C\|\phi\|_\infty.
\]
For every $v\in L^4$ and either cutoff interaction,
$\int|V^{(\cdot)}||v|^2d\mu\le\|V^{(\cdot)}\|_2\|v\|_4^2<\infty$.
Thus both vectors also lie in both multiplication-form domains.

\LS{2}{2}{The form Duhamel identity is}
\[
 \ip\phi{(T_K(s_2)-T_{K'}(s_2))\psi}
 =\int_0^{s_2}
 \ip{T_K(u)\phi}{\Delta V,T_{K'}(s_2-u)\psi}\,du.
\tag{5.11}
\]

\Because Differentiate
$u\mapsto\ip{T_K(u)\phi}{T_{K'}(s_2-u)\psi}$ in the common form domain.
The two $H_0$ form terms cancel, leaving multiplication by $\Delta V$.
Integrability follows from the estimates in the next step, so the fundamental
theorem for absolutely continuous form matrix elements applies.

\LS{2}{3}{For $0\le u\le s_4$, use exponents $(2,4,4)$ on
$(T_K(u)\phi,\Delta V,T_{K'}(s_2-u)\psi)$; for
$s_4\le u\le s_2$, reverse the two semigroup roles.  Then}
\[
 \|T_K(s_2)-T_{K'}(s_2)\|_{2\to2}
 \le18s_2M_0e^{2c_\lambda s_2}\|V^{(K)}-V^{(K')}\|_2.
\tag{5.12}
\]

\Because The long semigroup factor has elapsed time at least $s_4$ and is
$L^2\to L^4$ bounded by (5.10); the short factor is $L^2\to L^2$ bounded.
The interaction difference satisfies
$\|\Delta V\|_4\le9\|\Delta V\|_2$ by (1.13).  Integrate the common bound
over a length $s_2$ interval; the harmless factor $2$ accommodates both halves.

\LS{2}{4}{Equation (1.12) implies}
\[
 \|T_K(s_2)-T(s_2)\|_{2\to2}\to0.
\tag{5.13}
\]

\LS{2}{5}{This proves $\langle1\rangle12$.}\QedStep

\LS{1}{13}{Prove discreteness, eigenvalue convergence, trace class, and trace-norm
convergence for every positive time.}

Let $\nu_0\le\nu_1\le\cdots$ be the occupation-number eigenvalues of $H_0$.

\LS{2}{1}{$T(s_2)$ is compact, and if $\mu_j^{(K)}$ and $\mu_j$ are ordered
eigenvalues of $K^{(K)}$ and $K_L$, then}
\[
 \mu_j^{(K)}\to\mu_j,\qquad
 \mu_j\ge\tfrac12\nu_j-c_{1/2}.
\tag{5.14}
\]

\Because Each $T_K(s_2)$ is trace class by (3.7), hence compact.  The
operator-norm limit (5.13) is compact.  The min--max/singular-value inequality
gives convergence of every ordered heat eigenvalue
$e^{-s_2\mu_j^{(K)}}\to e^{-s_2\mu_j}$; the logarithm is continuous on
$(0,\infty)$.  Apply min--max to (2.22) at $\delta=1/2$ and pass to the
limit.

\LS{2}{2}{For every $s>0$,}
\[
 \Tr T_K(s)\to\Tr T(s)<\infty.
\tag{5.15}
\]

\Because By (5.14), each summand converges, while
\[
 e^{-s\mu_j^{(K)}}\le e^{sc_{1/2}}e^{-s\nu_j/2},\qquad
 \sum_je^{-s\nu_j/2}=Z_0(s/2)<\infty.
\]
Dominated convergence on the counting measure gives (5.15), and shows that the
spectrum of $K_L$ is discrete and exhausts the $\mu_j$.

\LS{2}{3}{For $u=s/2$,}
\[
 \|T_K(u)-T(u)\|_{\Jtwo}^2
 =\Tr T_K(2u)-2\Tr(T_K(u)T(u))+\Tr T(2u).
\tag{5.16}
\]

\Because Expand the square of the self-adjoint difference and use cyclicity for
products of a trace-class and a bounded operator.

\LS{2}{4}{The mixed trace in (5.16) converges to $\Tr T(2u)$.}

\Because If $(\phi_n)$ diagonalizes $T(u)$ with eigenvalues $\sigma_n$, then
\[
 \Tr(T_K(u)T(u))=\sum_n\sigma_n\ip{\phi_n}{T_K(u)\phi_n}.
\]
Each matrix element converges strongly, is bounded by $e^{uc_\lambda}$, and
$\sum_n\sigma_n=\Tr T(u)<\infty$; dominated convergence applies.

\LS{2}{5}{Thus $\|T_K(u)-T(u)\|_{\Jtwo}\to0$, and}
\[
 \|T_K(s)-T(s)\|_{\Jone}
 \le(\|T_K(u)\|_{\Jtwo}+\|T(u)\|_{\Jtwo})
 \|T_K(u)-T(u)\|_{\Jtwo}\to0.
\tag{5.17}
\]

\Because Factor
$T_K(s)-T(s)=T_K(u)[T_K(u)-T(u)]+[T_K(u)-T(u)]T(u)$ and apply the
Schatten H\"older inequality.  The Hilbert--Schmidt norms are bounded by (5.15).

\LS{2}{6}{For every bounded $B$,}
\[
 \Tr(T_K(s)B)\to\Tr(T(s)B).
\tag{5.18}
\]

\Because The absolute difference is at most
$\|T_K(s)-T(s)\|_{\Jone}\|B\|$.

\LS{2}{7}{This proves $\langle1\rangle13$.}\QedStep

\section{Periodic Feynman--Kac clause and Weyl insertions}

Fix a coarse Weyl symbol $\xi=\eps_Mq+ip$ and its continuum refinement.  Let
$H_p$ be the V2 continuum harmonic interpolant, $H_p^{(K)}$ its mode truncation,
and let $J_q,J_q^{(K)}$ be the corresponding sharp-time sources.

\LS{1}{14}{Pass every deterministic shift datum to the continuum cutoff limit.}

\LS{2}{1}{The action tail satisfies}
\[
 0\le S_\infty(H_p)-S^{(K)}(H_p)
 \le C_\xi(K/m)^{-\eta}.
\tag{6.1}
\]

\Because The tail is
\[
 \frac1{8L}\sum_{|k|>K}\gamma(k)\cth(t\gamma(k)/2)|b_\infty(k)|^2.
\]
The D4 envelope is $|b_\infty(k)|^2\le
C(1+\eps_M|k|)^{-2-\eta}$; $\cth(t\gamma/2)\le\cth(t_0m/2)$ and
$\gamma(k)\le m+|k|$.  Integral comparison of the resulting
$|k|^{-1-\eta}$ tail yields (6.1).

\LS{2}{2}{The source-phase tail satisfies}
\[
 \|\Phi(J_q)-\Phi(J_q^{(K)})\|_2^2
 \le C_\xi(K/m)^{-\eta},\qquad
 \|e^{i\Phi(J_q)}-e^{i\Phi(J_q^{(K)})}\|_2
 \le C_\xi(K/m)^{-\eta/2}.
\tag{6.2}
\]

\Because Its variance is the sharp-time thermal covariance paired with the source
tail:
\[
 \sum_{|k|>K}\frac{\cth(t\gamma(k)/2)}{2\gamma(k)}
 |a_\infty(k)|^2.
\]
Use the D4 envelope and integral comparison.  The second assertion uses
$|e^{ix}-e^{iy}|\le|x-y|$.

\LS{2}{3}{The shifted action difference obeys}
\[
 \|\U_t^{(K)}(\eta+H_p^{(K)})-
 \U_t(\eta+H_p)\|_2
 \le C_\xi[1+\log(1+K/m)]^2(K/m)^{-1/2}.
\tag{6.3}
\]

\LS{3}{1}{Expand every shifted Wick power using (3.12) and its quadratic
analogue.}

\LS{3}{2}{Terms in which the random Wick chaos is truncated are bounded by
(2.17)--(2.19).}

\LS{3}{3}{For deterministic coefficients,}
\[
 \sup_s\frac1{2L}\sum_{|k|>K}|\widehat H_p(s,k)|^2
 \le C_\xi(\eps_MK)^{-1-\eta},\qquad
 \|H_p^{(K)}\|_\infty,\|H_p\|_\infty\le C_\xi.
\tag{6.4}
\]

\Because The harmonic time factor has modulus at most one; square the D4 mode
envelope and sum its tail.  The uniform bound is the absolutely summable
$|k|^{-1-\eta/2}$ envelope.

\LS{3}{4}{The identity}
\[
 A^{*n}-B^{*n}=\sum_{r=0}^{n-1}A^{*r}*(A-B)*B^{*(n-1-r)}
\tag{6.5}
\]
and Young/Cauchy--Schwarz bounds turn (6.4) into the required differences of all
coefficient power symbols.

\Because Only $n\le4$ occurs and the untruncated coefficient sequences have
uniform $\ell^1$ bounds.  Each term of (6.5) contains exactly one tail factor,
so its $\ell^2$ norm has the rate in (6.4).  Wick isometry pairs it with the
bounded covariance convolution from V2.  The slower common rate in (6.3) is
therefore valid.

\LS{2}{4}{For every finite $p$,}
\[
 \sup_K\|e^{-\U_t^{(K)}(\eta+H_p^{(K)})}\|_p<\infty,\qquad
 e^{-\U_t(\eta+H_p)}\in L^p.
\tag{6.6}
\]

\Because Apply (1.4) pointwise with
$u=\eta^{(K)}(s,x)+H_p^{(K)}(s,x)$ and integrate: this gives
$\U_t^{(K)}(\eta+H_p^{(K)})\ge-Ct\ell_K^2$, uniformly in the admissible coarse
data.  Equation (6.3) is the required $L^2$ increment bound; use
$1+\log(1+K/m)\le C_\varsigma(K/m)^\varsigma$ with
$0<\varsigma<1/4$ to retain a positive power.  Step $\langle1\rangle5$ then gives
the uniform $L^p$ bound, and an almost-everywhere convergent subsequence followed by
Fatou gives the limiting assertion.

\LS{2}{5}{This proves $\langle1\rangle14$.}\QedStep

\LS{1}{15}{Prove the continuum periodic partition and marked-Weyl formulas.}

\LS{2}{1}{For $t\ge t_0$,}
\[
 \Tr e^{-tK_L}=Z_0(t)\E_{\mu^{\rm per}_{t,L}}[e^{-\U_t}]
 \in(0,\infty).
\tag{6.7}
\]

\Because Equation (3.8) holds at every $K$.  Its left side converges by (5.15).
For the right side, (4.2), (2.16), H\"older, and the uniform exponential bounds
give $e^{-\U_t^{(K)}}\to e^{-\U_t}$ in $L^1$.  Finiteness follows.  Positivity
follows because the exponential is positive a.s.; alternatively Jensen gives
$\E e^{-\U_t}\ge e^{-\E\U_t}>0$, and the finite chaos expectation is finite.

\LS{2}{2}{For the full refined coarse symbol,}
\[
 \frac{\Tr(e^{-tK_L}W_{\rm ct}(R_M^\infty\xi))}{\Tr e^{-tK_L}}
 =e^{-S_\infty(H_p)}
 \frac{\E[e^{i\Phi(J_q)}e^{-\U_t(\eta+H_p)}]}
 {\E[e^{-\U_t}]}.
\tag{6.8}
\]

\LS{3}{1}{Apply (3.10) to the symbol and interaction truncated at $K$.}

\LS{3}{2}{The Euclidean numerator converges.}

\Because Split its difference into a phase difference and an exponential
difference.  Bound the phase part by (6.2) times the uniform $L^2$ norm in
(6.6).  Bound the exponential part using (4.2), (6.3), and (6.6).  Equations
(6.1) and (6.7) handle the prefactor and denominator.

\LS{3}{3}{The operator numerator converges.}

\Because Split
\[
 \Tr[(T_K(t)-T(t))W(\xi_K)]
 +\Tr[T(t)(W(\xi_K)-W(R_M^\infty\xi))].
\]
The first term tends to zero by (5.17) and unitarity.  The one-particle symbols
$\xi_K$ converge in norm by the D4 envelope.  The Weyl relation and the Fock
vacuum overlap
\[
 \ip{W(\alpha)\Omega_0}{W(\beta)\Omega_0}
 =e^{i\sigma(\alpha,\beta)/2}e^{-\|\alpha-\beta\|^2/4}
\]
give strong continuity of $W(\xi_K)$.  Expanding the trace-class $T(t)$ in an
eigenbasis and using the bound $2e^{-t\mu_j}$ gives convergence of the second
term by dominated convergence.  Thus (6.8) follows.

\LS{2}{3}{This proves $\langle1\rangle15$.}\QedStep

\section{Positivity improvement, simple ground state, and positive gap}

\LS{1}{16}{Prove the continuum ground-state assertions.}

\LS{2}{1}{$T(t)$ is positivity preserving.}

\Because Each $T_K(t)$ has a nonnegative Feynman--Kac kernel; strong limits preserve
the closed positive cone in $L^2$.

\LS{2}{2}{If $f,g\ge0$ are nonzero and bounded, then}
\[
 \ip f{T(t)g}>0.
\tag{7.1}
\]

\Because Equation (4.5) has a nonnegative integrand.  The weight is strictly
positive a.s. because $\mathcal W$ is finite a.s. ($\mathcal W\in L^2$).  If
$A=\{f>0\}$ and $B=\{g>0\}$, then
\[
 \PP(X_0\in A,X_t\in B)
 =\ip{\mathbf1_A}{\Gamma(e^{-t\gamma})\mathbf1_B}>0
\]
by (2.3) and Simon's positivity-improving theorem, since
$\|e^{-t\gamma}\|=e^{-tm}<1$.  The integrand is strictly positive on an event
of positive probability.

\LS{2}{3}{Equation (7.1) extends to all nonzero $f,g\ge0$ in $L^2$.}

\Because Replace each by its bounded truncations $f\wedge n,g\wedge n$ and use
positivity preservation plus $L^2$ convergence.

\LS{2}{4}{The largest eigenvalue of the compact positive operator $T(t)$ is
simple and has an a.e. strictly positive eigenvector $\Omega_L$.}

\Because If $u$ is a top eigenvector, positivity gives
$T(t)|u|\ge|T(t)u|=\|T(t)\||u|$.  The variational characterization forces
equality, so $|u|$ is a nonnegative top eigenvector; (7.1) makes it strictly
positive.  If a second real top eigenvector $h$ were orthogonal to it, then both
$h_+$ and $h_-$ would be nonzero.  Positivity improvement gives
$T h_+>0$ and $T h_->0$ a.e., hence
$|Th|=|Th_+-Th_-|<Th_++Th_-=T|h|$ a.e.  Pairing with $\Omega_L>0$
contradicts equality of their top Rayleigh quotients.

\LS{2}{5}{If $\mu_0<\mu_1\le\cdots$ are the discrete eigenvalues of $K_L$,}
\[
 \gamma_L:=\mu_1-\mu_0>0.
\tag{7.2}
\]

\Because Simplicity says the eigenspace at $\mu_0$ has dimension one;
discreteness from Step $\langle1\rangle13$ prevents another spectral point from
accumulating at $\mu_0$.  Hence the next ordered eigenvalue is strictly larger.

\LS{2}{6}{The vector functional}
\[
 \omega^{\rm ct}_{L,\lambda}(A)=\ip{\Omega_L}{A\Omega_L}
\tag{7.3}
\]
is a state on the continuum Weyl algebra.

\Because Every Weyl operator whose symbol lies in
$H^{-1/2}(\T_L)\oplus iH^{1/2}(\T_L)$ is unitary on the Fock/Gaussian
representation.  A normalized vector functional is positive and unital; it extends
by continuity to the generated $C^*$-algebra.  The D4 refined coarse symbols lie in
this one-particle space by their decay envelope.

\LS{2}{7}{This proves $\langle1\rangle16$.}\QedStep

\section{V3 conclusion and scope fence}

\begin{resultbox}
Steps $\langle1\rangle9$--$\langle1\rangle10$ construct the bounded-below
self-adjoint $K_L$ and identify its operator and quadratic form on $\D$.
Steps $\langle1\rangle11$--$\langle1\rangle13$ prove trace class for every
positive time, trace-norm cutoff convergence, discrete spectrum, and eigenvalue
convergence.  Steps $\langle1\rangle14$--$\langle1\rangle15$ prove the periodic
partition and marked-Weyl formulas.  Step $\langle1\rangle16$ proves a simple
strictly positive ground vector and $\gamma_L>0$.

No assertion of essential self-adjointness of $(H_0+V_L)|_{\D}$, and no assertion
that $\D$ is a form core, is made or needed.  The source V3's compressed periodic
N2 tail count is repaired explicitly in (2.17)--(2.19).
\end{resultbox}

\end{document}
